% Propietario conservado: /Users/ruben/Documents/New project/output/REVISION_ARGUMENTAL_APERTURAS_CIERRES_20260915/VII/delivery/MANUSCRIPT_VII_ES_EN_COMPLETE/source_es/manuscrito/30c_composicion_corriente_conexion.tex
% Artículo VII, incorporación aditiva, revisión 05.
% Arquitectura autoral: energía de memoria APP--TRIT--TPK, S01.
% Antecedente recuperado: diferencial y refinamiento de 30b.
% Formalización reunida: nota COMPOSICION_VARIACIONAL_TRIT_CONEXION.md.
% Se conserva íntegra 30b; su representante antihermítico es la
% restricción unitaria del diferencial que aquí se desarrolla.
% Esta sección evalúa la corriente de una acción y una realización
% declaradas. No selecciona retrospectivamente su campo ni su peso.

\section{Réalisation du courant en coordonnées de connexion}
\label{vii:vii:sec:composicion-corriente-conexion}

L'énergie de mémoire attribue à chaque route une réponse qui conserve ses incidences et l'ordre des transports. Pour l'insérer dans une action géométrique, ses composantes par rapport aux variations de connexion sont nécessaires. Ce passage s'obtient en différentiant le transport qui réalise déjà la route : le courant de connexion est la composition de cette différentielle avec le covecteur énergétique de la section précédente.

Les régimes elliptique, parabolique et hyperbolique du TRIT requièrent de
conserver la différentielle complète avant de la restreindre à une
représentation particulière. Les rotations unitaires constituent un
domaine de cette construction. L'élargissement suivant conserve également
les directions de dilatation et de transport admises par la réalisation
géométrique utilisée.

\subsection{Différentielle sur les transports inversibles}

On maintient les fibres, les champs et le poids de
\eqref{vii:vii:eq:energia-memoria-discreta}. Considérons maintenant des transports
inversibles et des familles différentiables avec
\(\dot U_a=A_aU_a\), où \(A_a\) appartient à l'espace tangent
de la représentation choisie. Nous écrivons
\[
 v_a=U_af_{s(a)},\qquad r_a=f_{t(a)}-v_a,\qquad
 q=\mathsf W_mr,\qquad E_m=\langle r,\mathsf W_mr\rangle.
\]

\begin{proposition}[Représentant complet de la différentielle]
\label{vii:vii:prop:diferencial-invertible}
À champ et poids fixes, on a
\begin{equation}
 \mathrm d_UE_m(A)=-2\operatorname{Re}\sum_a
       \langle q_a,A_av_a\rangle
 =\sum_a\operatorname{Re}\operatorname{Tr}(\mathcal G_a^*A_a),
 \qquad \mathcal G_a=-2q_av_a^*.
 \label{vii:vii:eq:gradiente-transporte-completo}
\end{equation}
Les projections antihermitienne et hermitienne de ce représentant sont
\begin{equation}
 \frac{\mathcal G_a-\mathcal G_a^*}{2}
   =v_aq_a^*-q_av_a^*=\mathcal J_a,
 \qquad
 \frac{\mathcal G_a+\mathcal G_a^*}{2}
   =-(q_av_a^*+v_aq_a^*).
 \label{vii:vii:eq:proyecciones-gradiente-memoria}
\end{equation}
La restriction aux générateurs antihermitiens retrouve exactement
\eqref{vii:vii:eq:representante-corriente-discreta}.
\end{proposition}

\begin{proof}
L'identité \(\dot r_a=-A_av_a\) et le caractère autoadjoint du poids donnent \(\dot E_m=2\operatorname{Re}\langle q,\dot r\rangle\). En outre,
\[
 \operatorname{Tr}(\mathcal G_a^*A_a)
 =-2\operatorname{Tr}(v_aq_a^*A_a)=-2q_a^*A_av_a.
\]
Les deux projections s'obtiennent en prenant l'adjoint. Pour le produit
réel de Hilbert--Schmidt, les sous-espaces hermitien et antihermitien sont
orthogonaux ; la partie hermitienne a donc une contraction nulle avec
un générateur antihermitien.
\end{proof}

La formule permet des poids qui couplent les arêtes : leur action se calcule dans
\(q=\mathsf W_mr\) avant la séparation des contributions de chaque
incidence. Si le champ ou le poids varient, on conserve l'expression complète
\eqref{vii:vii:eq:variacion-completa-memoria}, également valable pour les
transports inversibles considérés ici.

\subsection{Représentations quadratiques du TRIT}

Dans une réalisation matricielle bidimensionnelle de \(\mathbb A_\tau\), on peut prendre
\[
 J_{+}=\begin{pmatrix}0&1\\-1&0\end{pmatrix},\qquad
 J_{-}=\begin{pmatrix}0&1\\1&0\end{pmatrix},\qquad
 J_{0}=\begin{pmatrix}0&1\\0&0\end{pmatrix}.
\]
Par multiplication directe,
\[
 J_+^2=-I,\qquad J_-^2=I,\qquad J_0^2=0.
\]
Leurs exponentielles réalisent respectivement les régimes elliptique, hyperbolique et parabolique. Le premier est orthogonal pour le produit positif fixé. Le second préserve la forme \(B=\operatorname{diag}(-1,1)\), puisque \(J_-^{\mathsf T}B+BJ_-=0\). Le troisième représente l'algèbre des nombres duaux. Ces représentations fixent des exemples des types quadratiques ; leur utilisation dans un écran d'une autre dimension conserve l'application de représentation correspondante.

Un calcul distingue les composantes de la différentielle. Sur une arête,
prenons \(U=I\), \(f_s=(1,1)^{\mathsf T}\),
\(f_t=2f_s\) et \(\mathsf W=I\). Alors \(q=v=f_s\),
\(\mathcal J=0\), tandis que
\begin{equation}
 \mathrm dE(J_-)=-4,\qquad \mathrm dE(J_0)=-2.
 \label{vii:vii:eq:control-regimenes-corriente}
\end{equation}
Le représentant complet conserve les deux réponses. La distinction
identifie l'espace des variations avant d'effectuer la contraction
matérielle.

\subsection{Composition des routes et inversion orientée}

\begin{proposition}[Transport contragrédient de la différentielle]
\label{vii:vii:prop:pullback-invertible}
Pour \(U_\gamma=U_N\cdots U_1\), soient
\(B_k=U_N\cdots U_{k+1}\) et \(B_N=I\). On a
\[
 \dot U_\gamma U_\gamma^{-1}
   =\sum_kB_kA_kB_k^{-1}.
\]
Le représentant \(\mathcal G_\gamma\) dans la fibre finale induit
à l'incidence \(k\)
\begin{equation}
 \mathcal G_k=B_k^*\mathcal G_\gamma(B_k^{-1})^*.
 \label{vii:vii:eq:transporte-contragrediente}
\end{equation}
\end{proposition}

\begin{proof}
La règle du produit, suivie de la multiplication par
\(U_\gamma^{-1}\), démontre la première égalité. La cyclicité
de la trace donne
\[
 \operatorname{Re}\operatorname{Tr}
  (\mathcal G_\gamma^*B_kA_kB_k^{-1})
 =\operatorname{Re}\operatorname{Tr}
  ((B_k^*\mathcal G_\gamma(B_k^{-1})^*)^*A_k).
\]
Pour \(B_k\) unitaire, on retrouve la conjugaison de
\ref{vii:vii:prop:composicion-respuesta}.
\end{proof}

\begin{proposition}[Inversion par congruence énergétique]
\label{vii:vii:prop:inversion-congruencia}
Pour une énergie locale par arête, l'inversion est représentée par
\begin{equation}
 U_{\bar a}=U_a^{-1},\qquad
 r_{\bar a}=-U_a^{-1}r_a,\qquad
 W_{\bar a}=U_a^*W_aU_a.
 \label{vii:vii:eq:inversion-peso-general}
\end{equation}
Elle conserve l'énergie. Si \(W_a\) reste fixe,
\[
 A_{\bar a}=-U_a^{-1}A_aU_a,\qquad
 \dot W_{\bar a}=U_a^*(A_a^*W_a+W_aA_a)U_a.
\]
La variation complète de l'énergie inversée coïncide avec la variation originale.
\end{proposition}

\begin{proof}
La substitution dans \(r_{\bar a}^*W_{\bar a}r_{\bar a}\) produit \(r_a^*W_ar_a\). La dérivée de l'inverse est \(\dot U_a^{-1}=-U_a^{-1}A_a\) ; celle du poids s'obtient par la règle du produit. L'égalité des énergies vaut pour toute la famille, de sorte que leurs dérivées coïncident, y compris le terme \(\langle r_{\bar a},\dot W_{\bar a}r_{\bar a}\rangle\).
\end{proof}

La congruence conserve la positivité pour tout transport inversible.
Dans la représentation unitaire, elle se réduit à la règle d'inversion de
\ref{vii:vii:prop:inversion-respuesta}. La somme utilise une orientation
par arête, ou la moitié de la somme sur les deux orientations.

\subsection{Évaluation des composantes de courant}

Soit \(\theta=(\theta_b)\) une variation réelle des coordonnées
de connexion de la réalisation utilisée. La différentielle effective du
transport définit les opérateurs
\begin{equation}
 A_a(\theta)=\dot U_aU_a^{-1}
     =\sum_bB_{ab}\theta_b.
 \label{vii:vii:eq:diferencial-conexion-transporte}
\end{equation}
Pour une composition finie, leurs contributions s'obtiennent par
\eqref{vii:vii:eq:transporte-contragrediente}. Les opérateurs
\(B_{ab}\) sont les dérivées de la représentation de connexion qui
réalise la route.

\begin{theorem}[Courant de connexion de l'action de mémoire]
\label{vii:vii:thm:corriente-conexion-evaluada}
Pour un terme matériel \(S_{\rm mem}=\mathfrak a_m E_m\), avec \(\mathfrak a_m\) sa normalisation dans la droite d'action, à champ, poids et normalisation fixes, on a
\begin{equation}
 \delta S_{\rm mem}=-\frac12\sum_b\theta_b s_b,\qquad
 s_b=4\mathfrak a_m\operatorname{Re}
        \sum_a\langle q_a,B_{ab}v_a\rangle.
 \label{vii:vii:eq:componentes-corriente-conexion}
\end{equation}
Si l'appariement cellulaire entre les bases de connexion et de courant
est une matrice inversible \(M\), définie par
\(\delta S_{\rm mem}=-\tfrac12\theta^{\mathsf T}M\sigma\),
les coordonnées du courant sont
\begin{equation}
 \sigma=M^{-1}s.
 \label{vii:vii:eq:coordenadas-corriente-material}
\end{equation}
\end{theorem}

\begin{proof}
On substitue \eqref{vii:vii:eq:diferencial-conexion-transporte} dans
\eqref{vii:vii:eq:gradiente-transporte-completo}, on multiplie par
\(\mathfrak a_m\) et on regroupe les coefficients de
\(\theta_b\). La convention variationnelle \(-1/2\) fixe le
facteur \(4\). Comparer les deux expressions pour toute variation
\(\theta\) donne \(M\sigma=s\) ; l'inversibilité de
l'appariement démontre l'existence et l'unicité de ces coordonnées.
\end{proof}

La formule réalise explicitement la composition mentionnée dans \ref{vii:vii:subsec:realizacion-respuesta-discreta}. Pour évaluer chaque composante, on transporte le champ, on calcule le résidu, on applique le poids et on contracte avec la variation de connexion de la même réalisation. L'accouplement \(M\) conserve l'orientation et le volume des cellules. Un accouplement diagonal donne la division correspondante par leurs volumes ; les autres bases conservent la matrice complète.

Si la connexion intervient aussi dans \(f\), \(\mathsf W_m\)
ou \(\mathfrak a_m\), on applique
\eqref{vii:vii:eq:variacion-completa-memoria} et l'on ajoute
\(E_m\delta\mathfrak a_m\). Les contributions de tous les
termes matériels sont additionnées avant de composer la réponse de Cartan.

\begin{corollary}[Compatibilité du courant avec le raffinement]
\label{vii:vii:cor:refinamiento-corriente-conexion}
Pour les familles compatibles de
\ref{vii:vii:thm:refinamiento-respuesta}, soit \(P_m\) l'application fixe
qui envoie les variations de connexion du niveau \(m\) au niveau
\(m+1\). Si les actions de mémoire coïncident sur ces champs
et transports raffinés, leurs covecteurs satisfont
\[
 s_m=P_m^{\mathsf T}s_{m+1},\qquad
 M_m\sigma_m=P_m^{\mathsf T}M_{m+1}\sigma_{m+1}.
\]
\end{corollary}

\begin{proof}
Différencier l'égalité des actions dans la direction
\(P_m\theta\) donne
\(-\tfrac12\theta^{\mathsf T}s_m
 =-\tfrac12\theta^{\mathsf T}P_m^{\mathsf T}s_{m+1}\).
Le caractère arbitraire de \(\theta\) et les représentations cellulaires
des deux covecteurs prouvent les deux identités. La composition de
ces applications itère l'égalité à tout nombre fini de raffinements.
\end{proof}

\subsection{Insertion dans l'équation de Cartan}

La section~\ref{vii:vii:sec:corriente-respuesta-cartan} normalise la trois-forme matérielle au moyen de \(\delta_\omega S_m=-\tfrac12\int\delta\omega^{IJ}\wedge
\sigma_{IJ}\). Le théorème précédent fournit ses composantes cellulaires pour l'action de mémoire et l'accouplement déclarés. La réponse de Cartan utilise ce courant avec la cotétrade, le paramètre de Barbero--Immirzi et le couplage gravitationnel de la même réalisation.

La dépendance matérielle à l'égard de la contorsion détermine comment
cette équation est résolue. Pour une matière affine, le courant est indépendant
de la contorsion et l'on applique l'inversion linéaire et l'élimination
quadratique démontrées dans la section~\ref{vii:vii:sec:corriente-respuesta-cartan}.
Pour une action d'
holonomie, la différentielle est évaluée à la connexion actuelle et cette dépendance
est conservée lors de la résolution de l'équation. Par exemple, une arête avec
\(f_s=f_t=1\), un poids un et \(U(t)=e^{it}\) possède
\[
 E(t)=2-2\cos t,\qquad E'(t)=2\sin t.
\]
Le courant dépend ici de \(t\). La formule complète maintient
cette dépendance au lieu de la remplacer par sa valeur dans une connexion
particulière.

Enfin, la densité gravitationnelle s'obtient à partir de la variation
métrique de l'action effective, avec son champ et sa normalisation.
La représentation du courant, la contorsion qu'il produit et la
densité résultante sont les opérations successives de cette composition.
La loi de dilution ultérieure transporte l'amplitude ainsi évaluée dans
la section matérielle correspondante.
